% titel: Blatt 5, A4: Intervallhalbierungsverfahren
% Teil II: der durchgerechnete Weg.

\subsection*{Gegeben und gesucht}
$p(x)=x^5-2x-1$, Nullstelle in $[0,2]$. Teil (a): drei Halbierungen von
Hand; Teil (b): 20 Schritte per Python und Nullstelle auf 5 Nachkommastellen.

\subsection*{Vorbedingung}
\emph{Gedankenschritt:} Bevor man halbiert, prüft man den Vorzeichenwechsel.
$p$ ist ein Polynom, also stetig; $p(0)=0-0-1=-1<0$, $p(2)=32-4-1=27>0$.
Der Zwischenwertsatz sichert damit eine Nullstelle in $[0,2]$.

\subsection*{a) Drei Schritte von Hand}
\emph{Gedankenschritt:} Bei jedem Schritt: Mittelpunkt $m$, Vorzeichen von
$p(m)$ gegen $p(\text{links})$ vergleichen, das Teilintervall mit dem
Vorzeichenwechsel behalten.
\begin{center}
\begin{tabular}{c|c|c|c|c}
Schritt & Intervall & $m$ & $p(m)$ & neues Intervall\\\hline
1 & $[0,2]$ & $1$ & $1-2-1=-2<0$ & $[1,2]$\\
2 & $[1,2]$ & $1{,}5$ & $7{,}59-3-1=3{,}59>0$ & $[1,1{,}5]$\\
3 & $[1,1{,}5]$ & $1{,}25$ & $3{,}05-2{,}5-1=-0{,}45<0$ & $[1{,}25,1{,}5]$\\
\end{tabular}
\end{center}
Nach 3 Schritten liegt die Nullstelle in $[1{,}25\,;\,1{,}5]$.

\subsection*{b) 20 Schritte per Python}
\emph{Warum dieser Code:} Der Produkt-Test $p(a)\cdot p(m)<0$ ist das Herz
der Auswahl; je nach Ergebnis wird die rechte oder linke Grenze auf $m$
gesetzt und das Intervall halbiert.
\begin{verbatim}
def p(x): return x**5 - 2*x - 1
a, b = 0.0, 2.0
for _ in range(20):
    m = (a + b) / 2
    if p(a) * p(m) < 0:
        b = m
    else:
        a = m
print((a + b) / 2)      # -> 1.2906488013...
\end{verbatim}

\begin{ergebnis}
Nullstelle $x^*\approx\boxed{1{,}29065}$ (auf 5 Nachkommastellen).
\end{ergebnis}

\begin{kontrolle}
\textbf{Fehlerbreite:} Nach 20 Schritten ist das Intervall
$2/2^{20}\approx1{,}9\cdot10^{-6}$ lang; die 6. Nachkommastelle ist damit
die unsichere. Der Mittelwert $1{,}2906488$ gerundet gibt $1{,}29065$.
\checkmark

\textbf{Probe:} $p(1{,}29065)\approx0{,}0000$ -- nahe der Nullstelle. An
den Intervallgrenzen schlägt das Vorzeichen gerade in der Nähe der
Nullstelle um; konsistent mit einer einzigen Nullstelle in $[1,2]$.
\checkmark
\end{kontrolle}

\begin{merksatz}
Bisektion braucht nur Stetigkeit und einen Vorzeichenwechsel; jeder Schritt
halbiert das Intervall, nach $n$ Schritten ist die Breite
$(b-a)/2^n$. Im Code ist der Produkt-Test $p(a)p(m)<0$ und das saubere
Überschreiben der beiden Grenzen entscheidend.
\end{merksatz}
